%! Linear Regression — Unit 1, Lesson 1.5 — Guided Notes (Answer Key)
% Mirrors notes/main.tex exactly; every \blank{} becomes an \ans{} and every work block is
% authored byte-identically, so blank and key paginate alike.
% Gradual release (course-wide shape, 2026-09-01): the notes are the "I do / we do" block,
% 20 minutes, and the Individual Practice block that closes them is the "you do alone", 15 more.
% There is no group activity.
\documentclass[10pt]{article}
\usepackage{algebra2-article}
\usepackage{algebra2-key}   % pulls in -boxes; do NOT also load -boxes
% Coordinate grid: {xmin}{xmax}{ymin}{ymax}
\newcommand{\lgrid}[4]{%
  \draw[step=1, linegray!40, very thin] (#1,#3) grid (#2,#4);
  \draw[->, semithick, charcoal] (#1,0)--(#2,0) node[below right, font=\scriptsize]{$x$};
  \draw[->, semithick, charcoal] (0,#3)--(0,#4) node[left, font=\scriptsize]{$y$};}
\newcommand{\dpt}[2]{\fill[navy] (#1,#2) circle (2.6pt);}
% The free-throw data: hours practised (x) vs. free throws made out of 15 (y).
% Vocab answer for the key: bold forest term, then the filled definition on a write-line.
% The leading and trailing \par are required: \ansline ends with \dotfill but does NOT end the
% paragraph, so without them each term label gets pulled onto the previous answer's line.
\newcommand{\vocabans}[2]{%
  \par\noindent\textbf{\textcolor{forest}{#1:}}\\[1pt]\ansline{#2}\par}
\newcommand{\ftdata}{\dpt{1}{4}\dpt{2}{4}\dpt{3}{7}\dpt{4}{7}\dpt{5}{11}\dpt{6}{11}\dpt{7}{12}\dpt{8}{14}}

\begin{document}

\pageheader{Unit 1, Lesson 1.5}{Guided Notes --- Answer Key}

\vspace{0.10in}

\begin{objectivebox}
\textbf{By the end of this lesson, I will be able to\dots}
\begin{itemize}\setlength{\itemsep}{2pt}
  \item read a \emph{scatterplot} and describe its association: \emph{direction}, \emph{form}, \emph{strength};
  \item interpret the \emph{correlation coefficient} $r$ --- sign for direction, size for strength;
  \item read a \emph{line of best fit} $\hat y=ax+b$, interpret its slope and intercept, and \emph{predict};
  \item tell \emph{interpolation} from \emph{extrapolation}, and say why \emph{correlation is not causation}.
\end{itemize}
\end{objectivebox}

\vspace{0.05in}

\begin{vocabbox}
\small
Fill in each term as we build it together.
\par\vspace{2pt}   % \par is required: \vocabans's \noindent is a no-op mid-paragraph
\vocabans{Scatterplot (bivariate data)}{paired measurements $(x,y)$ plotted as points, to see how two variables relate}
\vocabans{Association: direction, form, strength}{does $y$ rise or fall with $x$; do the points follow a line; how tightly}
\vocabans{Correlation coefficient $r$}{a number in $[-1,1]$: sign is the direction, size $|r|$ is the strength of the linear fit}
\vocabans{Line of best fit (regression line) $\hat y=ax+b$}{the line from technology that best summarises the trend; $\hat y$ is a predicted value}
\vocabans{Interpolation vs.\ extrapolation}{predicting inside the data's $x$-range (safe) vs.\ outside it (the trend may break)}
\vocabans{Correlation is not causation (lurking variable)}{a strong $r$ means the variables move together; a hidden third variable may drive both}
\end{vocabbox}

\begin{hookbox}
A coach records, for each player, the \emph{hours practised} this week ($x$) and the \emph{free
throws made} out of $15$ ($y$). The points do \emph{not} lie on one line --- but they trend upward.
\begin{center}
\begin{minipage}[c]{0.40\linewidth}\centering
\begin{tikzpicture}[scale=0.22]
  \lgrid{-0.5}{9}{-0.5}{15.5}
  \ftdata
\end{tikzpicture}
\end{minipage}\hfill
\begin{minipage}[c]{0.56\linewidth}
\begin{itemize}\setlength{\itemsep}{2pt}
  \item As hours practised increase, free throws made tend to \ans{increase} (increase / decrease).
  \item Could a \emph{single line} pass through every point? \ans{no} \; Could one line
        summarise the \emph{trend}? \ans{yes}
  \item The player who practised $5$ hours made \ans{$11$}; the one who practised $6$ made
        \ans{$11$}. So more practice does not guarantee more baskets --- it \emph{tends} to.
\end{itemize}
\end{minipage}
\end{center}
No exact rule fits messy data --- but a \textbf{line of best fit} can summarise the trend, and
today we learn to read it, measure it, and predict with it.
\end{hookbox}

\boxguard[24]
\begin{notesbox}{1. Scatterplots \& Association: Direction, Form, Strength}
A \textbf{scatterplot} plots paired data $(x,y)$. Describe its association three ways:
\textbf{direction} (as $x$ rises, does $y$ rise or fall?), \textbf{form} (do the points follow a
\emph{line}?), and \textbf{strength} (how tightly?). Classify each plot.
\begin{center}
\begin{tikzpicture}[scale=0.30]
  \lgrid{-0.5}{7}{-0.5}{7}
  \dpt{1}{2}\dpt{2}{2}\dpt{3}{4}\dpt{4}{4}\dpt{5}{6}\dpt{6}{6}
  \node[font=\scriptsize\bfseries] at (3,-1.8){(a)};
\end{tikzpicture}
\hspace{0.4cm}
\begin{tikzpicture}[scale=0.30]
  \lgrid{-0.5}{7}{-0.5}{7}
  \dpt{1}{6}\dpt{2}{5}\dpt{3}{5}\dpt{4}{3}\dpt{5}{2}\dpt{6}{2}
  \node[font=\scriptsize\bfseries] at (3,-1.8){(b)};
\end{tikzpicture}
\hspace{0.4cm}
\begin{tikzpicture}[scale=0.30]
  \lgrid{-0.5}{7}{-0.5}{7}
  \dpt{1}{3}\dpt{2}{6}\dpt{3}{2}\dpt{4}{5}\dpt{5}{3}\dpt{6}{5}
  \node[font=\scriptsize\bfseries] at (3,-1.8){(c)};
\end{tikzpicture}
\end{center}
\begin{itemize}[leftmargin=1.5em, itemsep=2pt]
  \item Plot (a): direction \ans{positive}; form: the points \ans{do} (do / do not) follow
        a line; strength \ans{strong}.
  \item Plot (b): direction \ans{negative} --- as $x$ rises, $y$ \ans{falls}; strength
        \ans{strong}.
  \item Plot (c): direction \ans{none} --- there is no clear linear trend at all.
\end{itemize}
The free-throw data is direction \ans{positive}, form \ans{linear}, strength \ans{strong}.
\end{notesbox}

\boxguard[26]
\begin{notesbox}{2. The Correlation Coefficient $r$ --- Direction \emph{and} Strength in One Number}
Technology reports one number, the \textbf{correlation coefficient} $r$, for how well a
\emph{line} fits the points. It always lies between \ans{$-1$} and \ans{$1$}. Its
\textbf{sign} gives the \emph{direction}; its \textbf{size} $|r|$ gives the \emph{strength}.
\begin{center}
\begin{tikzpicture}[scale=1]
  \draw[semithick, charcoal, ->] (-3.6,0)--(3.6,0);
  \foreach \x/\lab in {-3/-1, -1.5/-0.5, 0/0, 1.5/0.5, 3/1}{\draw (\x,0.12)--(\x,-0.12) node[below, font=\scriptsize]{$\lab$};}
  \node[font=\scriptsize, forest] at (-3,0.45){strong negative};
  \node[font=\scriptsize, forest] at (0,0.45){no linear trend};
  \node[font=\scriptsize, forest] at (3,0.45){strong positive};
\end{tikzpicture}
\end{center}
\textbf{Match each $r$ to a plot from Section 1} (write a, b, or c): \quad
$r\approx0.90$: \ans{a} \qquad $r\approx-0.85$: \ans{b} \qquad $r\approx0.05$: \ans{c}
\begin{center}
\renewcommand{\arraystretch}{1.35}
\begin{tabularx}{\linewidth}{@{}X|X@{}}
\textbf{$r=0.9$} & \textbf{$r=-0.9$} \\
\hline
Direction \ans{positive}; \; strength \ans{strong} & Direction \ans{negative}; \; strength \ans{strong} \\
Points hug a line that \emph{rises} & Points hug a line that \emph{falls} \\
\end{tabularx}
\end{center}
\vspace{2pt}
\begin{tcolorbox}[colback=goldbg, colframe=goldacc, boxrule=0.9pt, arc=2mm,
    left=2.5mm, right=2.5mm, top=1.5mm, bottom=1.5mm]
\small \textbf{Careful --- the sign is direction only; it says nothing about strength.} A
classmate says ``$r=-0.9$ is a weak fit, because it is negative.'' Which is the \emph{stronger}
linear fit, $r=-0.9$ or $r=0.4$? \ans{$r=-0.9$} \; Why? \; $|{-0.9}|=$ \ans{$0.9$} and $|0.4|=$
\ans{$0.4$}. The points with $r=-0.9$ sit \emph{much} closer to their line --- it just slopes
down. Two different questions, two different parts of $r$.
\end{tcolorbox}
\end{notesbox}

\boxguard[12]
\begin{notesbox}{3. The Line of Best Fit --- and What Its Numbers Mean}
Technology also returns the \textbf{line of best fit} (the \emph{regression line}) $\hat y=ax+b$,
the one line that best summarises the trend; the hat means a \emph{predicted} value. For the
free-throw data:
\[ \hat y = 1.5x + 2, \qquad r = 0.97. \]
\begin{center}
\begin{minipage}[c]{0.40\linewidth}\centering
\begin{tikzpicture}[scale=0.19]
  \lgrid{-0.5}{9}{-0.5}{15.5}
  \draw[very thick, forest] (0,2)--(8.6,14.9);
  \ftdata
\end{tikzpicture}
\end{minipage}\hfill
\begin{minipage}[c]{0.56\linewidth}
\textbf{Interpret in context --- this is the point.}
\begin{itemize}[leftmargin=1.5em, itemsep=2pt]
  \item The \textbf{slope} $a=1.5$: each additional \ans{hour} of practice is associated
        with about \ans{$1.5$} more free throws made --- a \emph{rate}, with units.
  \item The \textbf{$y$-intercept} $b=2$: the predicted free throws for a player who practised
        \ans{$0$} hours --- a \emph{baseline}.
  \item $r=0.97$: a \ans{strong} (strong / weak) fit, because $|r|$ is close to \ans{$1$}.
\end{itemize}
\end{minipage}
\end{center}
\end{notesbox}

\boxguard[28]
\begin{notesbox}{4. Predicting --- and Its Limits}
To \textbf{predict}, substitute an $x$ into $\hat y=ax+b$. The data ran from $x=1$ to $x=8$ hours.
\begin{center}
\begin{minipage}[t]{0.47\linewidth}
\textbf{Inside the data: $x=6$ hours.}
\begin{work}
  \hat y &= 1.5(6) + 2 \\
         &= 9 + 2 \\
         &= 11
\end{work}
$6$ is inside $1$ to $8$: this is \ans{interpolation}. \\ Reasonable? \ans{yes}
\end{minipage}\hfill
\begin{minipage}[t]{0.47\linewidth}
\textbf{Outside the data: $x=12$ hours.}
\begin{work}
  \hat y &= 1.5(12) + 2 \\
         &= 18 + 2 \\
         &= 20
\end{work}
$12$ is outside $1$ to $8$: this is \ans{extrapolation}. \\ Reasonable? \ans{no}
\end{minipage}
\end{center}
\begin{tcolorbox}[colback=goldbg, colframe=goldacc, boxrule=0.9pt, arc=2mm,
    left=2.5mm, right=2.5mm, top=1.5mm, bottom=1.5mm]
\small \textbf{Careful --- the line does not know the story's limits.} Only \ans{$15$} shots
were attempted, so ``$20$ made'' is impossible. Inside the data the trend held; outside it, the
line keeps going but reality does not --- always check an \emph{extrapolation} against the story.
\end{tcolorbox}
\vspace{2pt}
\textbf{Correlation is not causation.} In summer, ice-cream sales and pool drownings both rise, so
their $r$ is \ans{positive} (positive / negative) --- but ice cream does not \emph{cause} drownings.
The real driver, \ans{hot weather}, is a \textbf{lurking variable} pushing both. A strong $r$ shows
two variables \ans{move} together, not that one \ans{causes} the other.
\end{notesbox}

\boxguard[24]
\begin{practicebox}
A study relates \textbf{daily minutes of exercise} ($x$) to \textbf{resting heart rate in beats
per minute} ($y$) for adults who exercise $0$ to $45$ minutes a day. Technology gives
$\hat y=-0.4x+78$ with $r=-0.88$.
\begin{enumerate}[leftmargin=1.7em, itemsep=3pt]
  \item The association is \ans{negative} (positive / negative) and \ans{strong} (strong / weak),
        because the sign of $r$ is \ans{negative} and $|r|=$ \ans{$0.88$} is close to $1$.
  \item Interpret the slope $-0.4$ in context, with units.
        \par\ansline{Each additional minute of daily exercise is associated with a resting heart rate}
        \ansline{about $0.4$ beats per minute \emph{lower} --- a rate, not a guarantee for any one person.}
  \item Predict the resting heart rate for $x=30$ minutes.
        \begin{work}
          \hat y &= -0.4(30) + 78 \\
                 &= -12 + 78 \\
                 &= 66
        \end{work}
        $30$ is inside $0$ to $45$, so this is \ans{interpolation}. \; Does a strong $r$ prove exercise
        \emph{causes} a lower heart rate? \ans{no}
\end{enumerate}
\end{practicebox}

\boxguard[22]
\begin{scenariobox}[Individual Practice --- On Your Own]{navy}
Three problems, quietly and on your own. Sign is direction, size is strength; every slope gets
its units; check each prediction against the data's range.
\begin{enumerate}[leftmargin=1.9em, itemsep=6pt]
  \item \textbf{Match $r$ to the plot.} Assign each value $r\in\{\,0.95,\ 0.10,\ -0.70\,\}$.
    \begin{center}
    \begin{tikzpicture}[scale=0.24]
      \lgrid{-0.5}{6}{-0.5}{7}
      \dpt{1}{1}\dpt{2}{2}\dpt{3}{4}\dpt{4}{5}\dpt{5}{6}
      \node[font=\scriptsize\bfseries] at (2.5,-1.8){(P)};
    \end{tikzpicture}
    \hspace{0.4cm}
    \begin{tikzpicture}[scale=0.24]
      \lgrid{-0.5}{6}{-0.5}{7}
      \dpt{1}{4}\dpt{2}{2}\dpt{3}{5}\dpt{4}{3}\dpt{5}{4}
      \node[font=\scriptsize\bfseries] at (2.5,-1.8){(Q)};
    \end{tikzpicture}
    \hspace{0.4cm}
    \begin{tikzpicture}[scale=0.24]
      \lgrid{-0.5}{6}{-0.5}{7}
      \dpt{1}{5}\dpt{2}{3}\dpt{3}{4}\dpt{4}{2}\dpt{5}{2}
      \node[font=\scriptsize\bfseries] at (2.5,-1.8){(R)};
    \end{tikzpicture}
    \end{center}
    (P) $r=$ \ans{$0.95$} \quad (Q) $r=$ \ans{$0.10$} \quad (R) $r=$ \ans{$-0.70$} \quad
    Which plot has the strongest linear fit? \ans{P}

  \item \textbf{Interpret the model.} For workers with $0$ to $15$ years of experience, relating
        \emph{years of experience} ($x$) to \emph{salary in thousands of dollars} ($y$), technology
        gives $\hat y=2.5x+38$ with $r=0.91$. \\[3pt]
    (a) The slope $2.5$ means each additional \ans{year} of experience is associated with about
        \$\ans{$2{,}500$} more salary. \\[3pt]
    (b) The intercept $38$ means the predicted salary with no experience is \$\ans{$38{,}000$}. \\[3pt]
    (c) Predict the salary at $x=10$ years:
    \begin{work}
      \hat y &= 2.5(10) + 38 \\
             &= 25 + 38 \\
             &= 63
    \end{work}
    That is \$\ans{$63{,}000$}, and it is \ans{interpolation} (interpolation / extrapolation). \; A
    prediction at $x=30$ years would be \ans{extrapolation}.

  \item \textbf{Where the line breaks.} A caf\'e records the daily high \emph{temperature} ($x$,
        in $^\circ$F) and \emph{cups of hot cocoa sold} ($y$) on days from $20^\circ$ to $60^\circ$.
        Technology gives $\hat y=-3x+260$ with $r=-0.93$. \\[3pt]
    (a) Direction \ans{negative}; strength \ans{strong}. The slope $-3$ means: each extra degree
        is associated with about \ans{$3$} \ans{fewer} (more / fewer) cups sold. \\[3pt]
    (b) Predict at $x=40^\circ$ and at $x=90^\circ$:
    \begin{center}
    \begin{minipage}[t]{0.47\linewidth}
    \begin{work}
      \hat y &= -3(40) + 260 \\
             &= -120 + 260 \\
             &= 140
    \end{work}
    \end{minipage}\hfill
    \begin{minipage}[t]{0.47\linewidth}
    \begin{work}
      \hat y &= -3(90) + 260 \\
             &= -270 + 260 \\
             &= -10
    \end{work}
    \end{minipage}
    \end{center}
    (c) The $40^\circ$ prediction is \ans{interpolation}; the $90^\circ$ one is \ans{extrapolation}. Why is
        the second one unreasonable, and what should the caf\'e conclude instead?
    \par\ansline{You cannot sell $-10$ cups; $90^\circ$ is far outside the $20$--$60^\circ$ data, so the line has broken.}
    \ansline{Conclude only that very hot days sell few cups --- near zero --- not a negative number.}
\end{enumerate}
\end{scenariobox}

\end{document}
